\section{Limitations}\label{sec:limits}

\paragraph{The theory.} Model~F is an idealization: it assumes that every turn misses, that requests complete in FIFO order and
that a finished turn rejoins at once. The engine runs about 28 requests concurrently at 300k (14--90 across our pools), vLLM preempts, and hits return as the backlog
falls toward the onset, so model~F describes the trap, not the transitions into and out of it. We do not prove that the engine
converges to model~F, and the theorem's limits are reached only geometrically: in the registered window, about 33--37 generations
after the burst, a transient of order 10\% remains for our workload (MIRROR's late window, model~F's post-wave lull, lies 20--40\% below its asymptote), so we compare with the exact fluid as well as with the limits. The premium is
proved positive for our workload at $\rho\ge1.04$; below that it rests on a floating-point bracket and a first-order expansion.

\paragraph{The validation.} The check on the 57 stuck runs is retrospective: the theory was derived after they were read, and its
scoring tolerances were set on them. Only the runs submitted after tag \texttt{icml-trap-pred} test it blind, and only in outcome:
the four runs trapped by a tool slowdown (three TOOLSLOW-W, one TOOLSLOW-D15), the six CUT runs and MIRROR's runs, all on the A100
with the 300k pool. In the CUT runs each arm's registered rule rejected the naive reading, and CUT-A's also the quasi-static one;
what singles out model~F against intermediate readings is post hoc; MIRROR's rule rejects only the throughput-curve readings.
The from-empty runs are single runs per rate, and at 300k the
seed-0 bracket did not replicate on fresh seeds (\S\ref{sec:stampede}). In an exploratory replication on two plans chosen because
their outcomes were known, three runs each reproduced both originals, one collapse and one stable run, which suggests that the spread
across seeds at the edge comes from the plans rather than from service noise (six runs, on two GPUs of one node). The trapped windows
are at most 7500\,s long, and ``does not escape'' means within that horizon.

\paragraph{The experiments.} Sessions are synthetic (random tokens, exponential 5\,s tool times, geometric turn counts, fixed tool
outputs and output lengths, no shared prefixes), clients never retry or abandon, and real agent traces have tool times far shorter
and more dispersed~\citep{zhu2026tracelab}. Engines are colocated, with one GPU each, dense 7--8B models and pools of 150k to 956k tokens; disaggregated
serving, offloading, larger models and parallelism are untested. Seeds are shared across tests. The rescue engages a fixed cap at a
fixed time, the arrival stop ran at one rate, and the tool slowdowns at two. Appendix~\ref{app:threats} gives the full list.

\section{Conclusion}\label{sec:conclusion}

A burst can tip a prefix-cached agent server into a trap in which its cache stops paying; Lu et al., who posted first, and we
observed this independently (in their disaggregated system the hit rate falls to about 0.3, in ours every turn misses). Idealized as
model~F, the trap has an exact law: a linear renewal recursion along FIFO orbits carries every trapped state to the known
self-similar regime~\citep{khanin2001steady}, fixed by a Perron root, in which the gap between turns grows linearly, the backlog grows
at a load-dependent multiple of the excess load and the turns served are young and short; for our workload the multiple is 0.6 to 0.8 at our loads
and sessions complete faster than the all-miss capacity; with a 20-step budget they complete slower, and a blind test on the engine found completions below that capacity in the transient after a burst. Checked
retrospectively on 57 stuck runs on two engines, two GPUs and two models, with no parameter fitted to any trap measurement, these laws
hold for the runs taken together, and, with predictions frozen before the runs, on four runs in two cells trapped by a tool slowdown
alone, in four of the five pre-set scores (a per-run slope check failed in two). The practical reading is
Corollary~\ref{cor:trapthr}: in model~F, for our workload, a cut to just below the throughput measured in an all-miss trap does not clear it, and a cut
below the all-miss capacity drains the trap only slowly and by turning sessions away. On the engine, a cut between the two left all three CUT-A traps in place
to the end of the run, and a cut below the all-miss capacity shrank the three CUT-B backlogs only slowly, all still trapped at 8700\,s; each arm's registered rule rejected the naive reading, and what singles out
model~F against intermediate readings is post hoc. The other exit, a session cap that restores reuse, rescued all 9 established traps
of our registered tests at the same load without shedding, and the engine stayed good after the cap was removed.
